\chapter{Asincronismo e richieste di rete}

Completiamo la panoramica su JavaScript con tre argomenti: la
\textbf{memorizzazione dei dati} nel browser, i costrutti per la
\textbf{programmazione asincrona} e le \textbf{richieste di rete}.

\section{Memorizzare dati nel browser}

I browser moderni offrono diverse API utilizzabili da JavaScript per
memorizzare e recuperare i dati di un'applicazione: i \term{cookie}, il
\term{Web Storage} (\code{localStorage} e \code{sessionStorage}) e
\term{IndexedDB}.

\subsection{Cookie}

\dfn{Cookie}{
  I \term{cookie} sono piccole stringhe di dati. Fanno parte del protocollo HTTP
  e servono a \guillemotleft superarne\guillemotright\ la mancanza di stato. Vengono tipicamente
  impostati in una risposta HTTP tramite l'header \code{Set-Cookie}; i browser
  li memorizzano e li inviano in \textbf{tutte} le richieste successive verso lo
  stesso dominio, tramite l'header \code{Cookie}.}

È il meccanismo che avevamo anticipato nel primo capitolo, parlando della
natura \textit{stateless} di HTTP.

I cookie del sito corrente sono accessibili tramite la proprietà
\code{document.cookie}. Il suo valore è una stringa di coppie
\code{nome=valore} separate da \code{;}, e ogni coppia è un cookie distinto.

\begin{lstlisting}[style=html, name=cookie.html]
<h1>Hello!</h1>
<script>
  if (document.cookie.indexOf("name=") === -1) {   // nessun cookie "name"
    let name = prompt("Please, state your name:");
    document.cookie = `name=${name}; max-age=10`;  // scade fra 10 secondi
  }

  let name = document.cookie.replace("name=", "");
  document.querySelector("h1").innerHTML = `Hello, ${name}!`;
</script>
\end{lstlisting}

\subsection{localStorage}

Lavorare con i cookie da JavaScript non è comodissimo: soprattutto quando i
cookie sono molti, bisogna manipolare stringhe, spezzarle, usare espressioni
regolari.

L'oggetto \code{localStorage} offre un modo molto più semplice di memorizzare
coppie chiave/valore nel browser:

\begin{itemize}
  \item i dati \textbf{non} vengono inviati a ogni richiesta, quindi se ne
        possono memorizzare molti di più;
  \item l'interfaccia è simile a quella di una \code{Map};
  \item può memorizzare \textbf{solo valori stringa};
  \item esiste un \code{localStorage} distinto per ogni combinazione di dominio,
        protocollo e porta.
\end{itemize}

\begin{lstlisting}[style=html, name=localstorage.html]
<h1>Hello!</h1>
<script>
  if (localStorage.getItem("name") === null) {   // niente "name" memorizzato
    let name = prompt("Please, state your name:");
    localStorage.setItem("name", name);
  }

  let name = localStorage.name;
  document.querySelector("h1").innerHTML = `Hello, ${name}!`;
</script>
\end{lstlisting}

\subsection{sessionStorage e IndexedDB}

I dati in \code{localStorage} sopravvivono persino a un riavvio completo del
browser. \code{sessionStorage} offre le stesse funzionalità, ma si usa molto più
di rado: esiste solo all'interno della \textbf{scheda} corrente del browser, e i
dati sopravvivono a un aggiornamento della pagina ma non alla chiusura e
riapertura della scheda.

\term{IndexedDB} è invece un vero database integrato nel browser: supporta tipi
di chiave diversi, transazioni e interrogazioni su intervalli di chiavi, e può
memorizzare ancora più dati di \code{localStorage}. È inutilmente potente per le
tradizionali applicazioni web client-server, ed è pensato soprattutto per le
applicazioni che devono funzionare \textbf{offline}. Basta sapere che esiste.

\warning{Nulla di ciò che sta nel browser è sicuro}{
  Cookie, \code{localStorage} e IndexedDB vivono sul computer dell'utente e sono
  ispezionabili e modificabili con due clic dai DevTools. Non vanno quindi
  \textbf{mai} usati per memorizzare segreti, né per conservare informazioni di
  cui il server debba fidarsi: un carrello con i prezzi, per esempio, va sempre
  ricalcolato lato server. Vale la stessa regola vista per l'input dell'utente:
  ciò che arriva dal client va sempre verificato.}

\section{Sincronismo e asincronismo}

Il codice JavaScript viene tipicamente eseguito in modo \term{sincrono}: ogni
istruzione attende il completamento delle precedenti. È un modo di programmare
intuitivo, ma ha qualche inconveniente: alcune istruzioni possono richiedere
molto tempo, per esempio quando si attende l'input dell'utente, il
completamento di una richiesta di rete o la fine di un calcolo complesso. Le
istruzioni successive, magari importanti, restano bloccate ad aspettare.

\begin{lstlisting}[style=js, name=sincrono.js]
function firstOperation(value)  { return 1 + value; }
function secondOperation(value) { return 2 + value; }
function thirdOperation(value)  { return 3 + value; }

function setupPage() {
  let result = 0;
  result = firstOperation(result);
  result = secondOperation(result);
  result = thirdOperation(result);
  console.log(`Result is ${result}`);
}
\end{lstlisting}

\subsection{Callback}

\dfn{Callback}{
  Una \term{callback} è una funzione passata come argomento a un'altra
  funzione, con l'aspettativa che venga invocata al momento opportuno.}

Ci ricordano qualcosa? I gestori di eventi del capitolo precedente sono, in
effetti, delle callback.

Riscriviamo l'esempio precedente in stile asincrono, con le callback.

\begin{lstlisting}[style=js, name=callback.js]
function firstOperation(value, cb)  { cb(1 + value); }
function secondOperation(value, cb) { cb(2 + value); }
function thirdOperation(value, cb)  { cb(3 + value); }

function setupPage() {
  let result = 0;
  firstOperation(result, (r1) => {
    secondOperation(r1, (r2) => {
      thirdOperation(r2, (r3) => {
        console.log(`Result is ${r3}`);
      });
    });
  });
  console.log("Done");
}
\end{lstlisting}

\error{Callback Hell}{
  L'indentazione che si allarga a destra ha un nome: \term{Callback Hell},
  oppure \term{Pyramid of Doom}. Con molte callback annidate il codice diventa
  rapidamente illeggibile, e la situazione peggiora ulteriormente quando si deve
  gestire anche la propagazione degli errori: ogni livello deve controllare e
  inoltrare l'errore al livello successivo.

  Il JavaScript delle origini gestiva l'asincronismo così. Nel JavaScript
  moderno si usa un costrutto diverso: le \term{Promise}.}

\section{Promise}

Nella programmazione capita spesso di avere:

\begin{enumerate}
  \item del codice \term{produttore}, che fa qualcosa e ci mette un certo tempo:
        attendere l'input dell'utente, scaricare dati dalla rete, svolgere
        calcoli complessi;
  \item del codice \term{consumatore}, che ha bisogno dei risultati del
        produttore.
\end{enumerate}

Le \term{Promise} sono il modo di \guillemotleft collegare\guillemotright\ produttori e consumatori.

\subsection{Creare una Promise}

Un oggetto \code{Promise} si crea con il costruttore apposito.

\begin{lstlisting}[style=js]
let promise = new Promise(function (resolve, reject) {
  // Executor: qui il codice del produttore
});
\end{lstlisting}

La funzione passata come argomento si chiama \term{executor}, e viene invocata
\textbf{immediatamente} alla creazione della Promise. \code{resolve} e
\code{reject} sono callback fornite da JavaScript: quando il codice
dell'executor termina, deve invocare

\begin{itemize}
  \item \code{resolve(value)} se è andato a buon fine, con risultato
        \code{value};
  \item \code{reject(error)} se si è verificato un errore, dove \code{error} è
        l'oggetto errore.
\end{itemize}

\begin{figure}[H]
  \centering
  \begin{tikzpicture}[font=\Lato\small,
      s/.style={rounded corners=4pt, line width=1pt, align=center,
                minimum width=3.6cm, minimum height=1.15cm},
      a/.style={-{Stealth[length=6pt,width=5pt]}, line width=1pt}]

    \node[s, draw=neutral!55!white, fill=neutral!9!white] (p) at (0,0)
      {\scriptsize Status \textbf{"pending"}\\ \scriptsize Result
       \code{undefined}};

    \node[s, draw=green-gradient-6, fill=green-gradient-1!50!white] (f) at (6.6,1.5)
      {\scriptsize Status \textbf{"fulfilled"}\\ \scriptsize Result
       \code{result}};

    \node[s, draw=red-gradient-6, fill=red-gradient-1!40!white] (r) at (6.6,-1.5)
      {\scriptsize Status \textbf{"rejected"}\\ \scriptsize Reason
       \code{error}};

    \draw[a, draw=green-gradient-6] (p.east) to[out=20, in=180]
      node[wtlab, above, pos=0.55]{\code{resolve(result)}} (f.west);
    \draw[a, draw=red-gradient-6] (p.east) to[out=-20, in=180]
      node[wtlab, below, pos=0.55]{\code{reject(error)}} (r.west);

    \node[font=\Lato\scriptsize, text=neutral!35!black, anchor=south]
      at (0,0.75) {\code{new Promise(executor)}};
  \end{tikzpicture}
  \caption{Il ciclo di vita di una Promise.}
  \label{fig:promise-lifecycle}
\end{figure}

\info{Stato e risultato sono privati}{
  \textit{Status} e \textit{Result} sono proprietà \textbf{interne}: non si
  possono leggere direttamente dal codice. L'unico modo di accedere al risultato
  è registrare le funzioni da eseguire quando la Promise si risolve, come
  vediamo subito. Una Promise che ha lasciato lo stato \code{pending} si dice
  \term{settled} (\guillemotleft sistemata\guillemotright), e da quel momento il suo stato non
  cambia più.}

\begin{lstlisting}[style=js, name=promise.js]
let promise = new Promise(function (resolve, reject) {
  let number = Math.random();

  setTimeout(() => {     // setTimeout riceve una callback e un tempo in ms
    if (number > 0.5) {
      resolve(number);
    } else {
      reject(new Error("Number too small"));
    }
  }, 1000);
});

console.log(promise);    // Promise { <state>: "pending" }
\end{lstlisting}

\subsection{then, catch e finally}

Il codice consumatore riceve una Promise, cioè la promessa che un dato sarà
disponibile a un certo punto nel futuro. Le azioni da compiere quando la Promise
è risolta (o rifiutata) si registrano con i metodi offerti dall'oggetto Promise.
Il più importante è \code{.then()}.

\begin{lstlisting}[style=js]
promise.then(
  function (result) { /* invocata se la promise e' fulfilled */ },
  function (error)  { /* invocata se la promise e' rejected  */ }
)
\end{lstlisting}

Si può anche registrare una sola funzione, quella per il caso di successo.
Quando invece interessa solo il caso di errore si può passare \code{null} come
primo argomento, oppure — meglio — usare il metodo equivalente \code{.catch()}.

\begin{lstlisting}[style=js]
promise.catch(
  (error) => { /* invocata se la promise e' rejected */ }
)
\end{lstlisting}

Come i blocchi \code{try/catch}, anche le Promise hanno un metodo
\code{.finally()}. Riceve una callback senza argomenti, viene invocato appena la
Promise è \textit{settled} (in un senso o nell'altro), ed è pensato per le
operazioni di pulizia che vanno eseguite in ogni caso.

\begin{esempio}{L'ordine di esecuzione}
\begin{lstlisting}[style=js, name=promise-ordine.js]
let promise = new Promise(function (resolve, reject) {
  let number = Math.random();
  setTimeout(() => {
    if (number > 0.5) {
      resolve(number);
    } else {
      reject(new Error(`Number too small: ${number}`));
    }
  }, 1000);
});

console.log("Here");

promise.finally(() => { console.log("Promise settled") });

promise.then(
  (result) => { console.log(`Done: ${result}`) },
  (error)  => { console.log(error.message) }
);

console.log("There");
\end{lstlisting}

\begin{console}
Here
There
Promise settled
Done: 0.6      (oppure "Number too small: 0.3")
\end{console}

L'ordine è il punto centrale: \code{"There"} viene stampato \textbf{prima} del
risultato, perché il codice non si è fermato ad aspettare la Promise. È
esattamente ciò che volevamo ottenere con l'asincronismo.
\end{esempio}

\subsection{Concatenare le Promise}

E se dovessimo eseguire una \textbf{sequenza} di passi asincroni? È qui che
entra in gioco il \term{promise chaining}, e che si vede il vantaggio rispetto
alle callback annidate.

\begin{lstlisting}[style=js, name=chaining.js]
let promise = new Promise(function (resolve, reject) {
  setTimeout(() => { resolve(1); }, 1000);
});

console.log("Here");

promise
  .then((result) => { console.log(result); return result * 2; })
  .then((result) => { console.log(result); return result * 2; })
  .then((result) => { console.log(result); return result * 2; })

console.log("There");
\end{lstlisting}

\begin{console}
Here
There
1
2
4
\end{console}

Come funziona, se la Promise originale si risolve una sola volta con il valore
1? Il punto è che \textbf{ogni invocazione di \code{.then()} restituisce una
nuova Promise}, che rappresenta il completamento del gestore stesso. Quando un
gestore restituisce un valore, quel valore diventa il risultato della nuova
Promise, che alimenta il \code{.then()} successivo.

\begin{warningbox}{{Concatenare non è registrare più gestori}}
  I due frammenti seguenti sembrano simili ma fanno cose molto diverse.

\begin{lstlisting}[style=js]
// concatenazione: stampa 1, poi 2, poi 4
promise
  .then((r) => { console.log(r); return r * 2; })
  .then((r) => { console.log(r); return r * 2; })
  .then((r) => { console.log(r); return r * 2; });
\end{lstlisting}

\begin{lstlisting}[style=js]
// tre gestori indipendenti sulla STESSA promise: stampa 1, 1, 1
promise.then((r) => { console.log(r); return r * 2; });
promise.then((r) => { console.log(r); return r * 2; });
promise.then((r) => { console.log(r); return r * 2; });
\end{lstlisting}

  Nel primo caso si costruisce una catena, in cui ogni anello riceve il
  risultato del precedente. Nel secondo si registrano tre gestori indipendenti
  sulla \textbf{stessa} Promise, che ricevono tutti e tre lo stesso valore
  iniziale.
\end{warningbox}

\begin{esempio}{Errori in una catena}
\begin{lstlisting}[style=js, name=catena-errori.js]
promise.finally(() => {
  console.log("Promise settled");
}).then((r) => {
  console.log(r);
  return r * 2;
}).then((r) => {
  console.log(r);
  throw new Error("Oops");   // equivale a un reject()
  return r * 2;              // mai raggiunta
}).then((r) => {
  console.log(r);            // saltata: la catena e' in errore
  return r * 2;
}).catch((error) => {
  console.warn(error.message);
});
\end{lstlisting}

Quando un anello della catena fallisce, tutti i \code{.then()} successivi
vengono \textbf{saltati} fino al primo \code{.catch()}. È lo stesso
comportamento di un'eccezione che risale i blocchi \code{try}, e risolve
elegantemente il problema della gestione degli errori che rendeva pesanti le
callback.
\end{esempio}

\subsection{I metodi statici della classe Promise}

Spesso serve coordinare \textbf{più} operazioni asincrone. La classe
\code{Promise} offre quattro metodi statici per farlo.

\begin{center}
\renewcommand{\arraystretch}{1.4}
\begin{tabular}{@{}l@{\hspace{1.0cm}}>{\raggedright\arraybackslash}p{8.4cm}@{}}
  \textbf{Metodo} & \textbf{Si risolve\ldots} \\
  \code{Promise.all}        & quando \textbf{tutte} si risolvono; rifiuta appena
                              una rifiuta \\
  \code{Promise.allSettled} & quando tutte sono \textit{settled}, qualunque sia
                              l'esito \\
  \code{Promise.race}       & appena la \textbf{prima} si sistema, in un senso o
                              nell'altro \\
  \code{Promise.any}        & appena la prima si \textbf{risolve} (ignora i
                              rifiuti) \\
\end{tabular}
\end{center}

\begin{esempio}{Promise.all}
Prende in ingresso un array di Promise e restituisce una nuova Promise che
rappresenta il completamento di tutte. Il risultato è un array, in cui ogni
Promise di ingresso contribuisce con un elemento, \textbf{nell'ordine di
partenza} e non in quello di completamento.

\begin{lstlisting}[style=js, name=promise-all.js]
let prom = Promise.all([
  new Promise(resolve => setTimeout(() => resolve(3), 3000)),
  new Promise(resolve => setTimeout(() => resolve(2), 2000)),
  new Promise(resolve => setTimeout(() => resolve(1), 1000)),
]).then((result) => {
  console.log(result);   // Array(3) [3, 2, 1]
});
\end{lstlisting}
\end{esempio}

\begin{esempio}{Promise.allSettled}
\code{Promise.all} rifiuta se una qualunque delle Promise di ingresso rifiuta:
è un \textit{tutto o niente}. \code{Promise.allSettled} attende invece che tutte
si sistemino e restituisce un array di oggetti che descrivono lo stato di
ciascuna.

\begin{lstlisting}[style=js, name=promise-allsettled.js]
let p = Promise.allSettled([
  new Promise((res, rej) => { setTimeout(() => rej(new Error("Oops")), 1000) }),
  new Promise((res, rej) => { setTimeout(() => res(1), 2000) }),
]).then((result) => {
  console.log(result);
}).catch((err) => {
  console.warn(err.message);   // non viene eseguito
});
\end{lstlisting}
\end{esempio}

\begin{esempio}{Promise.race e Promise.any}
\code{Promise.race} si sistema appena una delle Promise di ingresso si sistema,
e ne assume il risultato (o l'errore).

\begin{lstlisting}[style=js, name=promise-race.js]
let p = Promise.race([
  new Promise((res, rej) => { setTimeout(() => rej(new Error("Oops")), 2000) }),
  new Promise((res, rej) => { setTimeout(() => res(1), 1000) }),
]).then((result) => {
  console.log(result);   // 1: la seconda arriva prima
});

let p2 = Promise.race([
  new Promise((res, rej) => { setTimeout(() => rej(new Error("Oops")), 1000) }),
  new Promise((res, rej) => { setTimeout(() => res(1), 2000) }),
]).catch((err) => {
  console.warn(err.message);   // Oops: qui arriva prima il rifiuto
});
\end{lstlisting}

\code{Promise.any} è simile, ma attende la prima Promise \textbf{risolta},
ignorando i rifiuti.

\begin{lstlisting}[style=js, name=promise-any.js]
let p = Promise.any([
  new Promise((res, rej) => { setTimeout(() => res(3), 3000) }),
  new Promise((res, rej) => { setTimeout(() => rej(new Error("Oops")), 1000) }),
  new Promise((res, rej) => { setTimeout(() => res(1), 2000) }),
]).then((result) => {
  console.log(result);   // 1
});
\end{lstlisting}

Se \textbf{tutte} rifiutano, \code{Promise.any} rifiuta con un
\code{AggregateError} che raccoglie tutti gli errori.

\begin{lstlisting}[style=js, name=promise-any2.js]
let p = Promise.any([
  new Promise((res, rej) => { setTimeout(() => rej(new Error("Woah")), 2000) }),
  new Promise((res, rej) => { setTimeout(() => rej(new Error("Oops")), 1000) }),
  new Promise((res, rej) => { setTimeout(() => rej(new Error("Ouch")), 1000) }),
]).catch((err) => {
  console.warn(err.message);             // No promise in Promise.any was resolved
  console.warn(err.errors[0].message);   // Woah
  console.warn(err.errors[1].message);   // Oops
  console.warn(err.errors[2].message);   // Ouch
});
\end{lstlisting}
\end{esempio}

Infine, \code{Promise.resolve(result)} e \code{Promise.reject(error)} creano
Promise \textbf{già sistemate}. Si usano soprattutto per compatibilità, quando
una funzione deve restituire una Promise anche se il valore è già disponibile.

\begin{lstlisting}[style=js]
let p = Promise.resolve(1);   // come new Promise(resolve => resolve(1))
p.then(result => console.log(result));
\end{lstlisting}

\section{async e await}

Il JavaScript moderno include una sintassi speciale per lavorare con le
Promise, basata su due parole chiave: \code{async} e \code{await}.

\subsection{async}

La parola chiave \code{async} si mette davanti a una funzione e garantisce che
quella funzione restituisca \textbf{sempre} una Promise:

\begin{itemize}
  \item qualunque altro valore restituito viene implicitamente avvolto in una
        Promise risolta;
  \item gli errori sollevati vengono avvolti in una Promise rifiutata.
\end{itemize}

\begin{lstlisting}[style=js, name=async.js]
async function f() {
  return new Promise(/* ... */);
}

async function g() {
  return 1;              // equivale a Promise.resolve(1)
}
\end{lstlisting}

\begin{lstlisting}[style=js, name=async-uso.js]
console.log("Here");

async function f() {
  // throw new Error("Ouch");
  return 1;
}

f().then(console.log).catch(console.warn);

console.log("There");
\end{lstlisting}

\begin{console}
Here
There
1
\end{console}

\subsection{await}

La parola chiave \code{await} si può usare \textbf{solo dentro funzioni async},
e sospende l'esecuzione finché una Promise non è sistemata.

\begin{lstlisting}[style=js, name=await.js]
console.log("Here");

async function f() {
  let promise = new Promise(function (res, rej) {
    setTimeout(() => { res("Done!") }, 2000);
  });

  let result = await promise;   // l'esecuzione si ferma qui
  console.log(`After Await. Result is ${result}`);
  return result;
}

f().then(console.log);

console.log("There");
\end{lstlisting}

\begin{console}
Here
There
After Await. Result is Done!
Done!
\end{console}

\info{Perché async/await è più leggibile}{
  \code{await} permette di scrivere codice asincrono che \textbf{si legge come
  codice sincrono}: niente catene di \code{.then()}, niente annidamenti. In più,
  dentro una funzione \code{async} si possono usare i normali blocchi
  \code{try/catch} per gestire gli errori, invece di \code{.catch()}.

  È importante però capire che si tratta di \textbf{zucchero sintattico} sopra
  le Promise: \code{await} non blocca il browser, sospende soltanto la funzione
  \code{async} in cui si trova, mentre il resto del programma continua a girare.
  Ecco perché nell'esempio \code{"There"} viene stampato prima del risultato.}

\section{Richieste di rete}

JavaScript può anche \textbf{scaricare dati} da Internet. Nel JavaScript
tradizionale si usavano oggetti \code{XMLHttpRequest} e delle callback per
gestire gli eventi di stato.

\begin{lstlisting}[style=html, name=xhr.html]
<h1>Cat Facts</h1>
<button id="btn">Click Here To Load a Cat Fact</button>
<p id="fact"></p>
<script>
  let btn = document.getElementById("btn");

  btn.onclick = () => {
    let req = new XMLHttpRequest();
    req.onload = handleFact;
    req.open("GET", "https://catfact.ninja/fact");
    req.send();

    function handleFact(result) {
      let fact = JSON.parse(this.responseText);
      document.getElementById("fact").innerHTML = fact.fact;
    }
  }
</script>
\end{lstlisting}

\subsection{L'API fetch}

Il modo moderno di effettuare richieste HTTP da JavaScript è la funzione
\code{fetch()}, che prende in ingresso una URL e un insieme opzionale di
opzioni, e restituisce una \textbf{Promise} che si risolve in un oggetto
\code{Response}.

\begin{lstlisting}[style=js]
let promise = fetch("url", [options]);
\end{lstlisting}

Le opzioni permettono di specificare il metodo HTTP, gli header della richiesta
e altro ancora: qui si ricongiungono le nozioni su HTTP viste nel primo
capitolo.

\subsection{Lavorare con la risposta}

La Promise restituita da \code{fetch} si risolve \textbf{appena arrivano gli
header} della risposta. A quel punto possiamo già controllare il codice di stato
e gli header, ma il \textbf{corpo} potrebbe non essere ancora disponibile.

\begin{lstlisting}[style=js, name=fetch-response.js]
async function f() {
  let response = await fetch("./example.json");

  // un header specifico
  console.log(response.headers.get('Content-Type'));   // application/json

  // iterazione su tutti gli header della risposta
  for (let [key, value] of response.headers) {
    console.log(`Header ${key}: ${value}`);
  }

  console.log(response.status);

  if (response.ok) {
    console.log("Request was successful");
  } else {
    console.log("Some error occurred");
  }
}
\end{lstlisting}

Per ottenere il corpo, l'oggetto \code{Response} offre diversi metodi, anch'essi
basati su Promise.

\begin{center}
\renewcommand{\arraystretch}{1.35}
\begin{tabular}{@{}l@{\hspace{1.1cm}}>{\raggedright\arraybackslash}p{8.0cm}@{}}
  \textbf{Metodo} & \textbf{Restituisce} \\
  \code{.text()}        & il corpo della risposta come testo \\
  \code{.json()}        & il corpo interpretato come JSON \\
  \code{.blob()}        & un \code{Blob} (dati binari con un tipo) \\
  \code{.arrayBuffer()} & un \code{ArrayBuffer} (rappresentazione binaria di
                          basso livello) \\
  \code{.formData()}    & un oggetto \code{FormData} (dati inviati da un form) \\
\end{tabular}
\end{center}

\error{Il corpo si può leggere una volta sola}{
  Si può scegliere \textbf{un solo} metodo di lettura del corpo. Se si è già
  invocato \code{response.text()}, una successiva \code{response.json()} non
  funzionerà, perché il contenuto del corpo è già stato consumato.}

\begin{esempio}{Leggere una risposta JSON}
\begin{lstlisting}[style=js, name=fetch-json.js]
async function f() {
  let response = await fetch("https://catfact.ninja/fact");

  if (response.ok) {
    let fact = await response.json();
    document.getElementById("fact").innerHTML = fact.fact;
  } else {
    console.log("Some error occurred");
  }
}
\end{lstlisting}

Da confrontare con la versione basata su \code{XMLHttpRequest} vista poco fa:
stesso risultato, molto meno codice e nessuna callback annidata.
\end{esempio}

\begin{esempio}{Scaricare dati binari con blob}
\begin{lstlisting}[style=html, name=fetch-blob.html]
<input id="msg" placeholder="Insert your message here" type="text">
<button id="btn">Generate QR</button>
<img id="img">
<script>
  document.getElementById("btn").onclick = () => {
    let msg = document.getElementById("msg").value;
    let url = `https://image-charts.com/chart?chs=200x200&cht=qr&chl=${msg}`;

    fetch(url)
      .then((response) => {
        return response.blob();
      })
      .then((blob) => {
        let img = document.getElementById("img");
        img.src = URL.createObjectURL(blob);
      });
  }
</script>
\end{lstlisting}

Il \code{Blob} scaricato viene trasformato in una URL locale temporanea con
\code{URL.createObjectURL()}, che si può assegnare direttamente all'attributo
\attr{src} di un'immagine.
\end{esempio}

\section{Riferimenti}

\begin{itemize}
  \item \textbf{The Modern JavaScript Tutorial} \\
        \urlx{https://javascript.info/} \\
        Parte 2: \textit{Storing data in the browser} (4.1, 4.2). \\
        Parte 1: \textit{Promises and async/await} (11.1--11.5, 11.8). \\
        Parte 3: \textit{Network requests} (3.1, 3.8).
  \item \textbf{Eloquent JavaScript} (3\textsuperscript{a} edizione), Marijn
        Haverbeke. \\
        \urlx{https://eloquentjavascript.net/} \\
        Capitoli 11 e 18.
  \item \textbf{JavaScript Reference} --- MDN web docs. \\
        \urlx{https://developer.mozilla.org/en-US/docs/Web/JavaScript/Reference}
\end{itemize}
